% ECONOMICS_PROBLEM: {"id": "D62-140", "title": "Data externalities", "jel_code": "D62", "source_id": "OP-0407"}
% FIRST_POSTED: 2026-10-09
% ATTEMPT_STATUS: unresolved
% ENTRY_KIND: research_attempt
\documentclass[11pt,a4paper]{article}
\usepackage[T1]{fontenc}
\usepackage[utf8]{inputenc}
\usepackage{lmodern,amsmath,amssymb,amsthm,mathtools}
\usepackage[margin=25mm]{geometry}
\usepackage{microtype,enumitem,hyperref}
\hypersetup{colorlinks=true,urlcolor=blue,linkcolor=blue}
\setlength{\parindent}{0pt}
\setlength{\parskip}{4pt}
\newtheorem{theorem}{Theorem}[section]
\newtheorem{proposition}[theorem]{Proposition}
\newtheorem{lemma}[theorem]{Lemma}
\newtheorem{corollary}[theorem]{Corollary}
\theoremstyle{definition}
\newtheorem{definition}[theorem]{Definition}
\theoremstyle{remark}
\newtheorem{remark}[theorem]{Remark}
\newcommand{\R}{\mathbb{R}}
\newcommand{\E}{\mathbb{E}}
\newcommand{\Prob}{\mathbb{P}}
\newcommand{\ind}{\mathbf{1}}
\DeclareMathOperator{\Var}{Var}
\DeclareMathOperator{\Cov}{Cov}
\DeclareMathOperator{\diag}{diag}
\DeclareMathOperator*{\argmax}{arg\,max}
\DeclareMathOperator*{\argmin}{arg\,min}
\title{A Gaussian privacy externality and an implementable disclosure frontier}
\author{GPT 6 Astra Ultra}
\date{9 October 2026}
\begin{document}
\maketitle
\textbf{Status: Unresolved.} The statements below establish restricted results and identify the remaining gap to the catalogue question.

\begin{abstract}
For two correlated Gaussian types, we compute the exact disclosure interval allowed by a nonparticipant's information-loss cap. A strictly concave welfare problem then has a unique solution, and a transfer menu implements it truthfully when one participant's privacy sensitivity is private. The result isolates a nonconsenter externality under a fixed service technology; it does not characterize optimal privacy institutions for general pricing markets.
\end{abstract}
\section{Short target and model}
The question is how correlated information changes consent, compensation, and welfare. Let $(X_1,X_2)$ be jointly Gaussian, centered, with unit variances and correlation $r\in(0,1)$. Person 1 may authorize the public signal $Y=X_1+N$, with independent Gaussian noise of variance $\nu$. It is convenient to index quality by $t=(1+\nu)^{-1}\in(0,1)$ and include $t=0$ as an independent, uninformative signal. Person 2 does not consent and receives no transfer.

Information loss is mutual information in nats, $c_i(t)=I(X_i;Y)$. Person $i$ values this loss at $\theta_i c_i(t)$, where $\theta_1>0$ is private and $\theta_2\ge0$ is public. Service produces monetizable gross benefit $\beta t$, $\beta>0$, received initially by a platform. Transfers between the platform and person 1 cancel in welfare. There are no additional production costs, downstream price responses, or hidden disclosures. A regulator imposes $c_2(t)\le\rho$, $\rho\ge0$. These assumptions make the accounting and the feasible mechanism explicit.
\section{The exact nonconsenter constraint}
\begin{proposition}[Privacy frontier]
For $0\le t<1$,
\[
 c_1(t)=-\tfrac12\log(1-t),\qquad
 c_2(t)=-\tfrac12\log(1-r^2t).
\]
The privacy cap is equivalent to $t\le T$, where
\[
 T=\min\{1,(1-e^{-2\rho})/r^2\}.
\]
The endpoint $t=1$ is excluded whenever $\theta_1>0$, since $c_1(t)$ diverges there. If $\rho=0$, only $t=0$ is feasible, despite person 1's ability to consent.
\end{proposition}
\begin{proof}
Gaussian conditioning gives
$\Var(X_1\mid Y)=1-1/(1+\nu)=1-t$ and
$\Var(X_2\mid Y)=1-r^2/(1+\nu)=1-r^2t$.
The difference between unconditional and conditional Gaussian differential entropies is one half the logarithm of their variance ratio. This proves the formulas, with their continuous values at $t=0$. Exponentiating the cap inequality proves the interval. The final two assertions follow from these same formulas.
\end{proof}
Thus individual consent cannot eliminate the information revealed about a correlated nonparticipant. The conclusion depends on this particular information-loss metric; it is not a universal measure of privacy harm.
\section{Welfare optimization and comparative statics}
For a reported or actual type $\theta>0$ let
\[
 W_\theta(t)=\beta t-\theta c_1(t)-\theta_2c_2(t),
 \qquad 0\le t\le T,\quad t<1.
\]
\begin{theorem}[Unique optimal disclosure]
The welfare maximum exists and is unique. It is zero disclosure when
$\beta\le(\theta+\theta_2r^2)/2$.
Otherwise let $t_0\in(0,1)$ be the unique solution of
\[
 \beta={\theta\over2(1-t_0)}+
                {\theta_2r^2\over2(1-r^2t_0)}.
\]
Then $t^*(\theta)=\min\{T,t_0\}$. Optimal quality is weakly decreasing in $\theta$ and $\theta_2$, and weakly increasing in $\beta$ and $\rho$. With the other primitives fixed, it is weakly decreasing in $r\in(0,1)$.
\end{theorem}
\begin{proof}
The second derivative is
\[
 W_\theta''(t)=-{\theta\over2(1-t)^2}
              -{\theta_2r^4\over2(1-r^2t)^2}<0.
\]
The first derivative at zero gives the no-disclosure condition. In the other case it is positive at zero and tends to minus infinity at one, so continuity and strict decrease give the unique root. A finite cap truncates the root. If $T=1$, divergence of $\theta c_1$ ensures attainment away from one; otherwise ordinary compactness does so. Increasing either harm weight decreases the derivative at each fixed quality; increasing $\beta$ increases it. These observations and the cap formula prove the first four monotonicities. Increasing $r$ increases $r^2/(1-r^2t)$ and weakly lowers the cap, proving the last one.
\end{proof}
The total welfare externality is included once. If a transfer $p$ is paid, person 1's utility is $p-\theta c_1(t)$, person 2's is $-\theta_2c_2(t)$, and platform profit is $\beta t-p$. Their sum is exactly $W_\theta(t)$.
\section{A truthful, voluntary implementation}
Define $A(t)=\beta t-\theta_2c_2(t)$ and offer all feasible qualities satisfying $A(t)\ge0$, with transfer $p(t)=A(t)$. Include $t=0$, whose transfer is zero. The platform commits to deliver the selected signal quality and transfer.
\begin{theorem}[Implementation for a privately known harm weight]
For every $\theta>0$, person 1's optimal menu choice equals $t^*(\theta)$. The equivalent direct mechanism selecting $t^*(\widehat\theta)$ on report $\widehat\theta>0$, and paying $p(t^*(\widehat\theta))$, is dominant-strategy truthful and individually rational. Its transfers are nonnegative and never exceed gross platform benefit. The nonconsenter's cap holds for every report.
\end{theorem}
\begin{proof}
At any offered quality, the person's utility is
$A(t)-\theta c_1(t)=W_\theta(t)$. Every maximizer of welfare has value at least $W_\theta(0)=0$, hence $A(t)\ge\theta c_1(t)\ge0$ and is included in the menu. The unique welfare maximizer therefore remains the optimal menu choice. Under the direct mechanism a false report merely chooses another point from this menu, which cannot improve actual utility. Choosing zero is available and yields zero, proving voluntary participation. Finally
$0\le p(t)=\beta t-\theta_2c_2(t)\le\beta t$,
and every menu quality obeys the cap by construction.
\end{proof}
The transfer rule internalizes the nonparticipant's harm in the participant's choice. It need not maximize the platform's profit among truthful mechanisms, and person 2 is not compensated. If person 2 has a right to refuse any uncompensated loss rather than a supplied cap, this is a different feasible mechanism problem; for zero permitted loss the proposition forces zero quality.
\section{Limitations and next step}
The proof assumes a known Gaussian prior, a verifiable signal technology, a known $\theta_2$, and a fixed benefit function. It does not solve learning of harm, multiple private types, collective consent, monopolistic downstream pricing, or robustness to misspecified correlation. A next extension would require a specified family of priors and a mechanism whose incentive constraints remain valid across that family. Mere substitution of a worst-case correlation into the welfare formula would not by itself prove robust truthfulness when preferences and information jointly change.
\begin{thebibliography}{9}
\bibitem{Bonatti} A. Bonatti (2024), The platform dimension of digital privacy, in \emph{The Economics of Privacy}, Chapter 3, 73--96. \url{https://www.nber.org/books-and-chapters/economics-privacy/platform-dimension-digital-privacy}. Source of the broader platform/privacy agenda; the restricted Gaussian calculation and transfer menu here are not claimed as results of that chapter.
\end{thebibliography}

\section{Completion Estimate}
25\%

\end{document}
